\scanpage{43}
d'où on conclut \unsure{(73)} $\nearrow$
et la formule à vérifier devient
\[
\begin{aligned}
\unsure{\mathcal N|\mathscr N|N}_{C_n/\mathcal O_S}(u1_C+vt)&=\deleted{\unsure{(-1)^n}}u^n-a_1u^{n-1}v\,\text{---}\,+\inserted{(-1)^n}a_nv^n\\
&=(-1)^nv^nf\left(\deleted{\unsure{+}}-\frac{u}{v}\right)
\end{aligned}\tag*{(\deleted{\unsure{73}}74)}
\]
Or si $u=0$, $v=1$ cette formule se réduit à
\[
 \mathcal N_{C_n/\mathcal O_S}(t)=\inserted{(-1)^n}a_n \tag*{(\deleted{\unsure{74}}75)}
\]
Or la matrice de la multiplication par $t$, dans
la base $1,t,\ldots,t^{n-1}$, \unsure{est}
\[
\begin{pmatrix}
0&0&&0&-a_n\\
1&0&&0&-a_{n-1}\\
0&1&&0&\\
\vdots&\vdots&&\vdots&\vdots\\
0&0&\cdots&0&\\
0&0&&1&-a_{\unsure{1|4}}
\end{pmatrix}
\]
dont le déterminant est $\underbrace{(-1)^{n+1}}_{\substack{\text{signature de la}\\\text{\deleted{\unsure{permuta}}permutation circulaire}\\\text{de l'ens à $n$ élts}}}(-a_n)=(-1)^{n\deleted{\unsure{+1}}}a_n$.

\deleted{Il suffit \unsure{pour} cela d'un} \inserted{On se ramène au cas de $t$ (au lieu de $u+vt$)}
par un argument universel. On peut aussi le
prouver directement, sans passer par des
déterminants, en faisant une ext. fid. plat
de la base permettant d'écrire
\[
 f\deleted{\unsure{_n}}(T)=\prod_{i=1}^n(T-t_i)\qquad\text{donc }a_i=(-1)^i\unsure{\tau|\sigma}_i(t_1,\ldots,t_n) \tag*{(76)}
\]
avec des sections $t_i$ de $\underline O_S$ $(1\leq i\leq n)$. \deleted{\unsure{On a}}
Alors (75) devient
\[
\begin{aligned}
\mathcal N_{C_n/\mathcal O_S}(u1_C+vt)&=(-v)^nf\left(\deleted{1}-\frac{u}{v}\right)\deleted{\unsure{v}}\;=\prod(u\deleted{\unsure{T}}+vt_i)\deleted{\unsure{(-1)^n}}\\
&=\prod_i \rho_i(u1_C+vt)
\end{aligned}\tag*{\deleted{\unsure{(77)}}}
\]

\noindent\rule{\linewidth}{0.3pt}

utilisant $\check V\xleftarrow{\sim} V\otimes L\simeq V$ \qquad $e_B\longmapsto T'$, $T\longrightarrow -e_B'$
(ici $(e_B',T'$ est la base duale de $e_B,T$)
% slanted
\margin{facteur de proportionnalité}
\[
 f_n=\underbrace{\varepsilon}_{\leftarrow}\deleted{\unsure{\lambda}}(-1)^n\left({e_B'}^n-a_1{e_B'}^{n-1}T'+\cdots+(-1)^ia_i{e_B'}^{n-i}{T'}^i\cdots+(-1)^na_n{T'}^n\right) \tag*{\deleted{\unsure{(73)}}}
\]
i.e.
\[
 f_n(ue_B+vT)=\deleted{\unsure{\varepsilon(-1)^n|\varepsilon}}\left(u^n-a_1u^{n-1}v+\cdots+(-1)^ia_iu^{n-i}v^i\,\text{---}\,+(-1)^na_nv^n\right) \tag*{(73)}
\]
(avec $\varepsilon=(-1)^n$ pour avoir $f_n(e_B)=1$)
% archive sheet number: 42
